\section*{Related Work}
\phantomsection\label{sec:related_work}

\paragraph{Execution Simulation and Verification.}
Execution trace generation improves code reasoning~\cite{nye2021scratchpad}, with extensions to output prediction~\cite{meta2025cwm,maimon2026selfexecution,liu2024codemind,liu2023codeexecutor,lee2026solidcoder}, failure analysis~\cite{abdollahi2025demystifying}, and error detection~\cite{li2024redo}. CRUXEval~\citep{gu2024cruxeval} benchmarks code execution prediction at function level; RepoExec~\citep{nam2025repoexec} evaluates repo-level code generation with execution-based correctness checks. Our work extends execution prediction to repo-level repair verification. Verification formulations span trajectory-level reward~\cite{shum2025swerm}, process-level feedback~\cite{li2025codeprm,xu2025patchreasoner}, per-test classification~\cite{ni2023lever,chen2023codet}, learned metrics~\cite{zhuo2024icescore,dong2024codescore,zhou2023codebertscore}, and RL with execution-based rewards~\cite{wang2026verpo,tang2026execverify}. Aletheia~\cite{venkatkrishna2026aletheia} compares training recipes; we compare \emph{formulations}. Test-time scaling~\cite{snell2024scaling,li2025sstar,chen2025vgsearch} and self-consistency~\cite{wang2023selfconsistency} demonstrate verification value; none compare formulations at matched capacity.

\paragraph{Process vs.\ Outcome Supervision.}
PRMs outperform ORMs for math reasoning~\cite{lightman2023letsverify,uesato2022solving}, with code-specific extensions~\cite{ye2025prlcoder,zhang2026funprm,jiang2025coderlplus,li2025codeprm}. GenRM~\cite{hosseini2025genrm} and GenPRM~\cite{zhao2026genprm} show generative verifiers outperform discriminative ones for math; our results reveal a domain difference---in code, the generative formulation does not outperform classification, which points to domain-specific tradeoffs (Appendix~\ref{app:threats}). In reward modeling terms, \cls{} maps to PRM, \traj{} to ORM, and \cwm{} to generative verification~\cite{hosseini2025genrm}; we also vary output format.

\paragraph{Repair, Mutation Testing, and Task Framing.}
CodeRL~\cite{le2022coderl}, AlphaCode~\cite{li2022alphacode,alphacode2023team}, and Reflexion~\cite{shinn2023reflexion} use execution feedback for repair. SWE-bench~\cite{jimenez2024swebench,chowdhury2024swebenchverified} benchmarks and SWE-agent~\cite{yang2024sweagent}/Agentless~\cite{xia2024agentless} demonstrate pipelines where verification is a bottleneck~\cite{yang2025aprsurvey}. Predictive mutation testing~\cite{tufano2019predictive,kim2022predictive,wang2024llmmutation,jain2023contextual} shares our per-test structure. \citet{raffel2020t5} showed that task framing affects NLP performance; we find that complexity moderates the classification--generation tradeoff in code verification.
